\noindent\textbf{(a)} $y=0$, y los arcos de $(y-x)^3=Cxy$, con $x\ne0$ y $x+2y\ne0$. El caso $C=0$ contiene $y=x$.

\noindent\textbf{(b)} $y=\pm x$ y las ramas $y=(Cx^2+C^{-1})/2$, $C\ne0$, restringidas por $Cx^2-C^{-1}\ge0$. Pueden incluirse empalmes $C^1$ con $y=\pm x$ en los puntos de contacto admisibles; sin esta posibilidad, la lista de soluciones es incompleta.
